
\subsection{分叉法与配边法：另一种证明}
\label{subsec:7-2-14}
\liftIndex{I328}

如第 5.3 节所述，
弗洛尔在 \textbf{[Flo88IV]} 中
证明弗洛尔同调不变性的原始论证
采用分叉论证；在 \textbf{[Flo89I]}
与 \textbf{[Oh93]} 中，
它被配边论证取代。
本书迄今用于定义滤过 $A_\infty$ 同态等对象的
方法，使用“时间有序积”
（\emph{top}，time ordered product）模空间，
因而更接近配边论证。
本小节说明一种以分叉论证为基础、
改用“逐时积”
（\emph{twp}，time-wise-product）模空间的替代方法。

考虑第 4.6 节或第~\ref{subsec:7-2-9} 小节的情形：
我们曾对辛微分同胚
$\psi\colon(M,L)\to(M',L')$
构造相应的滤过 $A_\infty$ 或
$A_{n,K}$ 同态。
其主要思路是使用“时间有序积”模空间
\[
\mathcal M^{\mathrm{main}}_{k+1}
(\{J_{\rho,s_0}\}_\rho:\beta;
\operatorname{top}(\rho))
\]
及由此模空间给出的对应关系，
构造滤过 $A_{n,K}$ 同态。

下面的替代方法使用“逐时积”模空间
$\mathcal M^{\mathrm{main}}_{k+1}
(\{J_\rho\}_\rho:\beta;
\operatorname{twp}(\rho))$，
其定义为
\begin{equation}
\begin{aligned}
&\mathcal M^{\mathrm{main}}_{k+1}
(\{J_\rho\}_\rho:\beta;
\operatorname{twp}(\rho))\\
&\qquad=\bigcup_{\rho\in[0,1]}
\{\rho\}\times
\mathcal M^{\mathrm{main}}_{k+1}
(M',L';J_\rho,\beta).
\end{aligned}
\tag{7.2.359}\label{eq:7-2-359}
\end{equation}
这里 $\{J_\rho\}_\rho$ 是 $M'$ 上的一族殆复结构，
满足 $J_0=\psi_*J$、$J_1=J'$。
以自然方式定义评价映射
\begin{equation}
\begin{aligned}
\operatorname{ev}
&=(\operatorname{ev}_1,\ldots,
\operatorname{ev}_k,\operatorname{ev}_0)\colon\\
&\mathcal M^{\mathrm{main}}_{k+1}
(\{J_\rho\}_\rho:\beta;
\operatorname{twp}(\rho))
\longrightarrow([0,1]\times L')^{k+1}.
\end{aligned}
\tag{7.2.360}\label{eq:7-2-360}
\end{equation}
仿照第~\ref{subsec:7-2-10} 小节，
取 $[0,1]\times L'$ 上光滑奇异单形的可数集
$\mathcal X_{\lambda,\lambda'}([0,1]\times L')$，
使每个 $\mathcal P$ 均满足
\begin{equation}
\psi_*^{-1}\mathcal P|_{\rho=0}
\in\mathcal X_{g(\lambda)}([0,1]\times L),
\tag{7.2.361.1}\label{eq:7-2-361-1}
\end{equation}
\begin{equation}
\mathcal P|_{\rho=1}
\in\mathcal X_{g(\lambda')}([0,1]\times L').
\tag{7.2.361.2}\label{eq:7-2-361-2}
\end{equation}
以上两个条件的靶集暂按源书保留
\translatornote{按 $\rho$ 端面识别，
$\mathcal P|_{\rho=0,1}$ 分别是
$L'$ 上的单形；经 $\psi^{-1}$ 后前者在 $L$ 上。
但原式仍把两处靶集写作带区间因子的
$\mathcal X_g([0,1]\times L)$、
$\mathcal X_g([0,1]\times L')$；
这与下文图式 7.2.35 的端点链模型不合，
此处不擅删区间因子。}。

再置
\begin{equation}
\begin{aligned}
&\mathcal M^{\mathrm{main}}_{k+1}
(\{J_\rho\}_\rho:\beta;
\operatorname{twp}(\rho);\vec{\mathcal P})\\
&\quad=\mathcal M^{\mathrm{main}}_{k+1}
(\{J_\rho\}_\rho:\beta;
\operatorname{twp}(\rho))
\mathbin{\times_{(\operatorname{ev}_1,\ldots,
\operatorname{ev}_k)}}
(\mathcal P_1\times\cdots\times\mathcal P_k).
\end{aligned}
\tag{7.2.362}\label{eq:7-2-362}
\end{equation}
令 $C([0,1]\times L';\lambda,\lambda')$
为由 $\mathcal X_{\lambda,\lambda'}
([0,1]\times L')$ 生成的自由
$\Lambda_{0,\mathrm{nov}}$ 模。
对 (7.2.362) 取适当扰动 $\mathfrak s$，
便在此模上得到滤过
$A_{n'(\lambda,\lambda'),K'(\lambda,\lambda')}$
代数结构。具体地，定义
\[
\mathfrak m_{k,\beta}
(\mathcal P_1,\ldots,\mathcal P_k)
=(\operatorname{ev}_0)_*\bigl(
\mathcal M^{\mathrm{main}}_{k+1}
(\{J_\rho\}_\rho:\beta;
\operatorname{twp}(\rho);\vec{\mathcal P})^{\mathfrak s}
\bigr).
\]
与前面讨论的情形一样，
须将右侧作单纯剖分，才能将其视为奇异链。
再利用 (7.2.361.1)、(7.2.361.2)
定义 $\operatorname{Eval}_{\rho=\rho_0}$
（$\rho_0=0,1$），得到图式~\ref{diag:7-2-35}：

\par\medskip\noindent\begin{minipage}{\textwidth}
\centering
$C(L,\lambda)
\xleftarrow{\ \operatorname{Eval}_{\rho=0}\ }
C([0,1]\times L';\lambda,\lambda')
\xrightarrow{\ \operatorname{Eval}_{\rho=1}\ }
C(L',\lambda')$
\par\smallskip
{\renewcommand{\thefigure}{7.2.35}\refstepcounter{figure}
\normalfont\bfseries 图式~\thefigure\label{diag:7-2-35}}
\end{minipage}\par\medskip

定义 $\operatorname{Eval}_{\rho=0}$ 时，
把 $\operatorname{ev}_0$ 与 $\psi^{-1}$ 复合。
$\mathfrak m_{k,\beta}$ 等的构造
与第~\ref{subsec:7-2-10} 小节的论证类似，
故省略细节。
因为 $\operatorname{Eval}_{\rho=0}$
是同伦等价（定理~\ref{thm:4-2-45}），
我们得到滤过
$A_{n'(\lambda,\lambda'),K'(\lambda,\lambda')}$
同态
\begin{equation}
C(L,\lambda)\longrightarrow C(L',\lambda').
\tag{7.2.363}\label{eq:7-2-363}
\end{equation}
这给出定理~\ref{thm:4-6-1} 的另一种证明。

严格说来，要从滤过
$A_{n'(\lambda,\lambda'),K'(\lambda,\lambda')}$
同态系构造滤过 $A_\infty$ 同态，
还必须证明不同 $\lambda,\lambda'$ 下的相容性。
由于定义 (7.2.363) 所用的滤过
$A_{n'(\lambda,\lambda'),K'(\lambda,\lambda')}$
同态依赖殆复结构族等辅助选择，
为构造这些滤过 $A_{n,K}$ 同态之间的同伦，
还需要同伦的同伦。
这可仿照第~\ref{subsec:7-2-12}、
\ref{subsec:7-2-13} 小节完成。

再往前一步，亦即为证明
定理~\ref{thm:4-6-25}，
需要研究滤过 $A_{n,K}$ 结构的
同伦的同伦的同伦。
这同样可仿照第~\ref{subsec:7-2-12}、
\ref{subsec:7-2-13} 小节完成。
